\section{Conclusion et suites à donner}

MEDICA-SERVICES dispose de fondations réelles, avec une stratégie approuvée, une PSSI, des sauvegardes planifiées et des équipes compétentes. Ses pratiques restent toutefois dépendantes des personnes et ne produisent pas les preuves qui permettraient à la direction de vérifier ses propres affirmations. L'écart de deux à trois niveaux avec la cible \Cible\ ne se comblera pas en une année. Le plan proposé traite d'abord les risques les plus graves et ramène les processus les plus faibles au niveau 2, dans le respect du plafond.

Quatre demandes de preuves conditionnent la confirmation de nos conclusions. Il s'agit du dossier d'investigation d'AL-1048 avec les journaux bruts, de l'export exhaustif des comptes privilégiés y compris SaaS, des procès-verbaux de restauration des sept services sans preuve de test, et des historiques de déploiement des correctifs avec la version installée sur APP01. Leur production pourra faire évoluer certains niveaux, notamment le niveau 2 de la stratégie aujourd'hui indéterminé, sans remettre en cause le niveau adéquat.
